% Lean: EconLean/Bridges/Aggregation/{GrossSubstitutes,LawOfDemand,Gorman,Examples}.lean
\chapter{Aggregation: Restrictions that Escape SMD (Bridge B0)}
\label{chap:aggregation}

The Sonnenschein--Mantel--Debreu theorem says that the standard consumer axioms restrict aggregate excess demand
only by continuity, homogeneity, Walras' law and boundary behaviour. This chapter formalizes micro restrictions that
(i) survive adding up consumers and (ii) buy an aggregate property (PLAN \S7.2a).

\section{Gross substitutes}
\label{sec:gross_substitutes}

\begin{definition}[Gross substitutes; MWG Def.~17.F.2]
\label{def:gross_substitutes}
\lean{Econ.PositivePrices, Econ.GrossSubstitutes, Econ.HomogeneousZero}
\leanok
$z$ has the \emph{gross substitutes} property if for all $p \gg 0$, all goods $l$ and all $t > p_l$, raising $p_l$ to
$t$ strictly raises $z_k$ for every $k \neq l$. $z$ is \emph{homogeneous of degree zero} if $z(tp) = z(p)$ for
$p \gg 0$, $t > 0$.
\end{definition}

\begin{theorem}[Gross substitutes is preserved under aggregation]
\label{thm:gs_sum}
\lean{Econ.GrossSubstitutes.add, Econ.GrossSubstitutes.sum}
\leanok
\uses{def:gross_substitutes}
The sum of two gross-substitute functions is gross substitute; so is $\sum_{a \in s} z_a$ for a nonempty finite set
$s$ of gross-substitute $z_a$.
\end{theorem}
\begin{proof}
\leanok
Add the strict inequalities (\texttt{linarith}; for a nonempty finite sum, \texttt{Finset.sum\_lt\_sum\_of\_nonempty}).
Nonemptiness is needed: the empty sum $0$ is not gross substitute.
\end{proof}

\begin{lemma}[Monotone path]
\label{lem:gs_path}
\lean{Econ.mixPrices, Econ.GrossSubstitutes.le_mixPrices, Econ.GrossSubstitutes.lt_of_le}
\leanok
\uses{def:gross_substitutes}
Let $z$ be gross substitute, $0 \ll p \le p'$, $p \ne p'$ and $p_k = p'_k$. Then $z_k(p) < z_k(p')$.
\end{lemma}
\begin{proof}
\leanok
For a finite set $S \not\ni k$ let $q_S$ equal $p'$ on $S$ and $p$ elsewhere. By induction on $S$
(\texttt{Finset.induction\_on}) $z_k(p) \le z_k(q_S)$, strictly if some $i \in S$ has $p_i < p'_i$: adding $l$ to $S$
changes only the $l$-th price, from $p_l$ to $p'_l$; if $p_l < p'_l$ gross substitutes gives a strict increase,
otherwise nothing changes. Apply this to $S = \iota \setminus \{k\}$, where $q_S = p'$ because $p_k = p'_k$, and some
$i \ne k$ has $p_i < p'_i$ because $p \neq p'$.
\end{proof}

\begin{theorem}[Uniqueness under gross substitutes; MWG Prop.~17.F.3]
\label{thm:gs_unique}
\lean{Econ.GrossSubstitutes.unique_equilibrium}
\leanok
\uses{def:gross_substitutes}
If $z$ is gross substitute and homogeneous of degree zero, and $z(p) = z(p') = 0$ with $p, p' \gg 0$, then
$p' = t p$ for some $t > 0$.
\end{theorem}
\begin{proof}
\leanok
\uses{lem:gs_path}
Choose $k$ maximizing $p_i/p'_i$ (\texttt{Finset.exists\_max\_image}) and let $\lambda = p_k / p'_k > 0$, $q = \lambda p'$.
Then $q \ge p$ with $q_k = p_k$, and $z(q) = z(p') = 0$ by homogeneity. If $q = p$ then $p' = \lambda^{-1} p$. Otherwise
Lemma~\ref{lem:gs_path} gives $0 = z_k(p) < z_k(q) = 0$, a contradiction.
\end{proof}

\section{The law of demand and the weak axiom}
\label{sec:law_of_demand}

\begin{definition}[Walras' law, WARP, law of demand; MWG Def.~17.F.1, \S4.C]
\label{def:law_of_demand}
\lean{Econ.WalrasLaw, Econ.ExcessWARP, Econ.LawOfDemand, Econ.WeakLawOfDemand}
\leanok
On a set of prices $D$: Walras' law is $p \cdot z(p) = 0$; WARP is: $z(p) \ne z(p')$ and $p \cdot z(p') \le 0$ imply
$p' \cdot z(p) > 0$; the (strict) law of demand is $(p' - p) \cdot (z(p') - z(p)) < 0$ whenever $z(p) \ne z(p')$; the
weak law has $\le 0$ for all $p, p'$.
\end{definition}

\begin{theorem}[Law of demand preserved under aggregation; MWG Prop.~4.C.1]
\label{thm:lod_sum}
\lean{Econ.LawOfDemand.weak, Econ.LawOfDemand.sum}
\leanok
\uses{def:law_of_demand}
If every $z_a$, $a \in s$, satisfies the strict law of demand on $D$, so does $\sum_{a\in s} z_a$.
\end{theorem}
\begin{proof}
\leanok
The strict law implies the weak one. If the aggregate differs at $p, p'$, some $z_b$ differs. By linearity
$(p'-p)\cdot(\sum_a z_a(p') - \sum_a z_a(p)) = \sum_a (p'-p)\cdot(z_a(p') - z_a(p))$; every term is $\le 0$ and the
$b$-term is $< 0$ (\texttt{Finset.sum\_lt\_sum}).
\end{proof}

\begin{lemma}[Law of demand implies WARP under Walras' law]
\label{lem:lod_warp}
\lean{Econ.LawOfDemand.excessWARP}
\leanok
\uses{def:law_of_demand}
With Walras' law on $D$, the strict law of demand implies WARP. (Caution: on a price cone with homogeneous $z$ these
hypotheses force $z \equiv 0$; the law of demand is meant for demand at fixed wealth or a normalized slice.)
\end{lemma}
\begin{proof}
\leanok
Expand $(p'-p)\cdot(z(p') - z(p)) < 0$ by bilinearity and use $p \cdot z(p) = p' \cdot z(p') = 0$:
$p' \cdot z(p) + p \cdot z(p') > 0$. If $p \cdot z(p') \le 0$ then $p' \cdot z(p) > 0$.
\end{proof}

\begin{theorem}[WARP makes the equilibrium set convex; MWG Prop.~17.F.1]
\label{thm:warp_convex}
\lean{Econ.ExcessWARP.convex_equilibria}
\leanok
\uses{def:law_of_demand}
If $D$ is convex and $z$ satisfies Walras' law and WARP on $D$, then $\{p \in D : z(p) = 0\}$ is convex.
\end{theorem}
\begin{proof}
\leanok
Let $z(p) = z(p') = 0$ and $q = ap + bp'$ ($a, b \ge 0$, $a + b = 1$), so $q \in D$. Suppose $z(q) \ne 0$. Walras' law
at $q$ gives $a\, p\cdot z(q) + b\, p' \cdot z(q) = 0$. If $p \cdot z(q) \le 0$, WARP (applied to $p, q$) gives
$q \cdot z(p) > 0$, impossible since $z(p) = 0$. Otherwise $p \cdot z(q) > 0$ forces $p' \cdot z(q) \le 0$ (else the sum
is positive; the case $a = 0$ is handled separately), and WARP applied to $p', q$ gives the same contradiction.
\end{proof}

\section{Gorman aggregation}
\label{sec:gorman}

Prices are fixed and suppressed; $x_a : \R \to \R^\iota$ is consumer $a$'s demand as a function of wealth.

\begin{definition}[Aggregation through wealth; Gorman form]
\label{def:gorman}
\lean{Econ.AggregatesThroughWealth, Econ.AggregatesThroughWealthNonneg, Econ.GormanForm}
\leanok
Demand \emph{aggregates through wealth} if there is $X$ with $\sum_a x_a(w_a) = X(\sum_a w_a)$ for all wealth
distributions $w$ (respectively, all $w \ge 0$). Demands have \emph{Gorman form} if $x_a(w) = x_a(0) + w \beta$ for a
common $\beta$.
\end{definition}

\begin{theorem}[Gorman form implies aggregation; MWG Prop.~4.B.1, ``if'']
\label{thm:gorman_if}
\lean{Econ.aggregatesThroughWealth_of_gorman}
\leanok
\uses{def:gorman}
Gorman-form demands aggregate through wealth, with $X(W) = \sum_a x_a(0) + W\beta$.
\end{theorem}
\begin{proof}
\leanok
Substitute and use $\sum_a w_a \beta = (\sum_a w_a)\beta$ (\texttt{Finset.sum\_smul}).
\end{proof}

\begin{lemma}[Transfer invariance]
\label{lem:gorman_transfer}
\lean{Econ.sum_update_apply, Econ.AggregatesThroughWealth.transfer,
  Econ.AggregatesThroughWealthNonneg.transfer}
\leanok
\uses{def:gorman}
If demand aggregates through wealth and $a \ne b$, then $x_a(u + h) + x_b(v - h) = x_a(u) + x_b(v)$ (for the
nonnegative version, whenever $u, v, u+h, v-h \ge 0$).
\end{lemma}
\begin{proof}
\leanok
Compare the distributions $w_0$ (giving $u$ to $a$, $v$ to $b$, $0$ to others) and $w_1$ (giving $u+h$ and $v-h$). They
have the same total, so the same aggregate demand; expanding both sums with the update identity leaves the claim.
\end{proof}

\begin{lemma}[Pairwise Gorman lemma]
\label{lem:gorman_pair}
\lean{Econ.AggregatesThroughWealth.pair, Econ.AggregatesThroughWealthNonneg.pair}
\leanok
\uses{lem:gorman_transfer}
If demand aggregates through wealth, demands are continuous (on $[0,\infty)$ in the nonnegative version) and $a \ne b$,
there is $\gamma$ with $x_a(u) = x_a(0) + u\gamma$ and $x_b(u) = x_b(0) + u \gamma$ (for $u \ge 0$ in the nonnegative
version).
\end{lemma}
\begin{proof}
\leanok
Put $v = h$ in Lemma~\ref{lem:gorman_transfer}: $x_a(u + h) = x_a(u) + g(h)$ with $g(h) = x_b(h) - x_b(0)$. Then $g$ is
additive: $g(h_1 + h_2) = g(h_1) + g(h_2)$, by comparing $x_a(h_1 + h_2)$ computed in two ways. $g$ is continuous, so by
the Cauchy equation (\texttt{map\_real\_smul} for a continuous additive map) $g(h) = h\, g(1)$. In the nonnegative
version, $g$ is only additive on $[0, \infty)$; extend it by $G(h) = g(h^+) - g(h^-)$, which is additive on $\R$ because
$(h_1+h_2)^+ + h_1^- + h_2^- = (h_1+h_2)^- + h_1^+ + h_2^+$ (from \texttt{max\_zero\_sub\_max\_neg\_zero\_eq\_self}) and
continuous, and agrees with $g$ on $[0,\infty)$. Take $\gamma = g(1)$.
\end{proof}

\begin{theorem}[Gorman's aggregation theorem; MWG Prop.~4.B.1, ``only if'']
\label{thm:gorman}
\lean{Econ.gorman_of_aggregatesThroughWealth, Econ.gorman_of_aggregatesThroughWealthNonneg}
\leanok
\uses{def:gorman}
With at least two consumers, if demand aggregates through wealth and each demand is continuous in wealth, demands have
Gorman form. The nonnegative version assumes aggregation and continuity only for $w \ge 0$ and concludes the Gorman form on
$[0, \infty)$. MWG assumes differentiability; continuity suffices here.
\end{theorem}
\begin{proof}
\leanok
\uses{lem:gorman_pair}
Fix $a_0$ and set $\beta = x_{a_0}(1) - x_{a_0}(0)$. For $b \ne a_0$, Lemma~\ref{lem:gorman_pair} gives a common slope
$\gamma$ for $a_0$ and $b$, and evaluating at $u = 1$ shows $\gamma = \beta$. For $a = a_0$ use any $b \ne a_0$
(\texttt{exists\_ne}); for $a \ne a_0$ use the pair $(a_0, a)$.
\end{proof}

\section{Examples}
\label{sec:aggregation_examples}

\begin{theorem}[Gross substitutes and WARP example]
\label{thm:zcd_example}
\lean{Econ.AggregationExamples.zCD, Econ.AggregationExamples.zCD_homogeneous, Econ.AggregationExamples.zCD_walras,
  Econ.AggregationExamples.zCD_grossSubstitutes, Econ.AggregationExamples.zCD_unique,
  Econ.AggregationExamples.zCD_warp}
\leanok
\uses{thm:gs_unique, thm:warp_convex}
$z(p) = (p_1/p_0 - 1,\, p_0/p_1 - 1)$ is homogeneous, satisfies Walras' law, gross substitutes and WARP on $p \gg 0$.
Every equilibrium is a multiple of $(1,1)$, and the equilibrium set is convex.
\end{theorem}
\begin{proof}
\leanok
Direct computation (\texttt{field\_simp}). For WARP, with $D = p_0p'_1 - p_1p'_0$:
$p\cdot z(p')\,p'_0p'_1 = (p'_1 - p'_0)D$ and $p'\cdot z(p)\,p_0p_1 = -(p_1 - p_0)D$; different excess demands force
$D \ne 0$, and the sign conditions follow by cases on the sign of $D$.
\end{proof}

\begin{theorem}[Law of demand, Gorman and production examples]
\label{thm:aggregation_examples}
\lean{Econ.AggregationExamples.lawOfDemand_neg, Econ.AggregationExamples.xG, Econ.AggregationExamples.xG_budget,
  Econ.AggregationExamples.xG_gorman, Econ.AggregationExamples.techY, Econ.AggregationExamples.convex_techY,
  Econ.AggregationExamples.prodEconomy, Econ.AggregationExamples.prodEconomy_eq,
  Econ.AggregationExamples.prodEconomy_paretoOptimal, Econ.AggregationExamples.uTwo,
  Econ.AggregationExamples.prodEconomy2, Econ.AggregationExamples.prodEconomy2_eq,
  Econ.AggregationExamples.rational_two, Econ.AggregationExamples.lns_two,
  Econ.AggregationExamples.prodEconomy2_paretoOptimal}
\leanok
\uses{thm:lod_sum, thm:gorman, thm:first_welfare_production, thm:second_welfare_production}
$z(p) = -p$ satisfies the strict law of demand. $x_a(w) = (a, -a) + w(\tfrac12, \tfrac12)$ is a budget-balanced
Gorman system at prices $(1,1)$. With technology $Y = \{y : y_0 \le 0,\ y_1 \le -y_0\}$: a one-consumer economy where
all hypotheses of both production welfare theorems hold; and a two-consumer economy ($u = x_0 + 2x_1$, endowment
$(1,0)$, half shares) whose equilibrium at $p = (1,1)$ has $y = (-2, 2)$, $x_a = (0,1)$, and is Pareto optimal.
\end{theorem}
\begin{proof}
\leanok
Direct computation. For $-p$: $(p'-p)\cdot(-p' + p) = -\|p' - p\|^2 < 0$. In the production examples, profit
$y_0 + y_1 \le 0$ on $Y$; on the budget line $y_0 + y_1 \le 1$, $u(y) \le 2(y_0 + y_1) \le 2 = u(0,1)$.
\end{proof}

\section{Tâtonnement stability}
\label{sec:tatonnement}

\begin{definition}[Tâtonnement path]
\label{def:tatonnement}
\lean{Econ.priceDistSq, Econ.IsTatonnementPath}
\leanok
A tâtonnement path for $z$ is $p : \R \to \R^\iota$ with $\dot p(t) = z(p(t))$ for all $t \ge 0$ (one-sided at $0$). $V(p) = \sum_i (p_i - p^*_i)^2$.
\end{definition}

\begin{lemma}[Lyapunov derivative; norm conservation]
\label{lem:tatonnement_deriv}
\lean{Econ.hasDerivWithinAt_priceDistSq, Econ.tatonnement_hasDerivWithinAt, Econ.tatonnement_deriv, Econ.tatonnement_norm_const}
\leanok
\uses{def:tatonnement}
If Walras' law $p(t) \cdot z(p(t)) = 0$ holds along a tâtonnement path, then
$\frac{d}{dt} V(p(t)) = -2\, p^* \cdot z(p(t))$ for $t \ge 0$, and $\sum_i p_i(t)^2 = \sum_i p_i(0)^2$.
\end{lemma}
\begin{proof}
\leanok
Componentwise, $\frac{d}{dt}(p_i - p^*_i)^2 = 2(p_i - p^*_i) z_i(p)$. Summing,
$\dot V = 2\, p \cdot z(p) - 2\, p^* \cdot z(p) = -2\, p^* \cdot z(p)$ by Walras' law (one-sided derivative on
$[0, \infty)$, a two-sided one for $t > 0$). With $p^* = 0$ the derivative is $0$; a function continuous on $[0,\infty)$
with zero derivative inside is both non-increasing and non-decreasing there, hence constant.
\end{proof}

\begin{theorem}[Tâtonnement stability, Lyapunov step; MWG Prop.~17.H.1]
\label{thm:tatonnement}
\lean{Econ.tatonnement_antitone, Econ.tatonnement_strictAnti}
\leanok
\uses{lem:tatonnement_deriv}
Suppose Walras' law holds along a tâtonnement path, and $p^* \cdot z(p(t)) > 0$ whenever $z(p(t)) \ne 0$. Then
$t \mapsto \|p(t) - p^*\|^2$ is non-increasing on $[0, \infty)$, and strictly decreasing if $z(p(t)) \ne 0$ for all
$t > 0$.
Convergence to the equilibrium ray is not claimed.
\end{theorem}
\begin{proof}
\leanok
By Lemma~\ref{lem:tatonnement_deriv} the function is continuous on $[0,\infty)$ with derivative $-2\, p^* \cdot z(p(t))$
for $t > 0$: zero at equilibrium points and negative elsewhere. By the mean value theorem it is antitone (strictly
antitone) on $[0, \infty)$.
\end{proof}

\begin{theorem}[Two goods: gross substitutes gives $p^* \cdot z(q) > 0$]
\label{thm:gs_cost_pos}
\lean{Econ.GrossSubstitutes.cost_pos_of_ratio_lt, Econ.GrossSubstitutes.cost_pos_two_goods}
\leanok
\uses{lem:gs_path}
Let $z$ be gross substitute, homogeneous of degree zero, with Walras' law on positive prices, and $p^* \gg 0$ an
equilibrium. With two goods, $p^* \cdot z(q) > 0$ at every $q \gg 0$ with $z(q) \ne 0$.
\end{theorem}
\begin{proof}
\leanok
Let $r_i = q_i / p^*_i$. If $r_0 = r_1 = c$ then $q = c p^*$ and $z(q) = z(p^*) = 0$ by homogeneity, a contradiction.
Otherwise say $r_k < r_j$ ($\{k, j\}$ the two goods) and let $q' = r_k p^*$, an equilibrium by homogeneity. Then
$q' \le q$, $q'_k = q_k$, $q'_j < q_j$, so Lemma~\ref{lem:gs_path} gives $z_k(q) > 0$. Walras' law gives
$z_j(q) = -q_k z_k(q) / q_j$, so
$p^* \cdot z(q) = p^*_k z_k(q)\,(1 - r_k p^*_j / q_j) = p^*_k z_k(q)(1 - r_k / r_j) > 0$.
\end{proof}

\begin{theorem}[Tâtonnement example]
\label{thm:zcd_tatonnement}
\lean{Econ.AggregationExamples.zCD_cost_pos, Econ.AggregationExamples.zCD_tatonnement_antitone,
  Econ.AggregationExamples.φArc, Econ.AggregationExamples.zArc, Econ.AggregationExamples.pArc,
  Econ.AggregationExamples.φArc_bounds, Econ.AggregationExamples.hasDerivAt_φArc,
  Econ.AggregationExamples.arctan_pArc, Econ.AggregationExamples.zArc_pArc,
  Econ.AggregationExamples.pArc_isTatonnementPath, Econ.AggregationExamples.pArc_walras,
  Econ.AggregationExamples.sin_lt_cos_φArc, Econ.AggregationExamples.pArc_cost_pos,
  Econ.AggregationExamples.zArc_pArc_ne, Econ.AggregationExamples.pArc_strictAnti}
\leanok
\uses{thm:gs_cost_pos, thm:tatonnement, thm:zcd_example}
(i) For $z(p) = (p_1/p_0 - 1, p_0/p_1 - 1)$, every tâtonnement path that stays in positive prices for $t \ge 0$ moves
weakly closer to $(1,1)$. (ii) A nonconstant witness: $z(q) = (\pi/4 - \arctan(q_1/q_0))(-q_1, q_0)$ satisfies Walras'
law, and $p(t) = (\cos\varphi(t), \sin\varphi(t))$ with $\varphi(t) = \pi/4 - (\pi/8)e^{-t}$ is a tâtonnement path that
is never at equilibrium and moves strictly closer to $(1,1)$.
\end{theorem}
\begin{proof}
\leanok
(i) Theorem~\ref{thm:gs_cost_pos} with $p^* = (1,1)$, using the gross substitutes, homogeneity and Walras' law of
Theorem~\ref{thm:zcd_example}, then Theorem~\ref{thm:tatonnement}. (ii) For $t \ge 0$, $\varphi(t) \in [\pi/8, \pi/4)$,
so $\arctan(\tan\varphi) = \varphi$ and $z(p(t)) = \tfrac{\pi}{8}e^{-t}(-\sin\varphi, \cos\varphi) = \dot p(t)$ by the chain
rule. Walras' law: $\cos\varphi(-\sin\varphi) + \sin\varphi\cos\varphi = 0$. And
$p^* \cdot z(p(t)) = \tfrac{\pi}{8}e^{-t}(\cos\varphi - \sin\varphi) > 0$ because $\sin\varphi < \sin(\pi/2 - \varphi) =
\cos\varphi$ for $\varphi < \pi/4$. Theorem~\ref{thm:tatonnement} applies.
\end{proof}

\section{Tâtonnement with any number of goods}
\label{sec:tatonnement_sup}

\begin{lemma}[Extreme relative prices]
\label{lem:gs_extreme_ratio}
\lean{Econ.GrossSubstitutes.excess_nonpos_of_max_ratio, Econ.GrossSubstitutes.excess_neg_of_max_ratio,
  Econ.GrossSubstitutes.excess_nonneg_of_min_ratio}
\leanok
\uses{lem:gs_path}
Let $z$ have gross substitutes and be homogeneous of degree zero, and let $p^* \gg 0$ with $z(p^*) = 0$. At any
$q \gg 0$, if $k$ maximizes $q_i/p^*_i$ then $z_k(q) \le 0$ (strictly if $q$ is not a multiple of $p^*$), and if $k$ minimizes it then $z_k(q) \ge 0$.
\end{lemma}
\begin{proof}
\leanok
Let $c = q_k/p^*_k$ and $q' = c p^*$, so $z(q') = 0$. In the max case $q \le q'$ with $q_k = q'_k$. If $q = q'$ then
$z_k(q) = 0$; otherwise Lemma~\ref{lem:gs_path} gives $z_k(q) < z_k(q') = 0$. The min case is symmetric.
\end{proof}

\begin{lemma}[Danskin-type monotonicity]
\label{lem:sup_antitone}
\lean{Econ.antitoneOn_sup'_of_hasDerivWithinAt}
\leanok
Let $f_1, \dots, f_n$ have right derivatives $f_i'(t)$ within $[0,\infty)$. If at every $t \ge 0$ each $f_i$ attaining
$\max_j f_j(t)$ has $f_i'(t) \le 0$, then $M(t) = \max_i f_i(t)$ is nonincreasing on $[0,\infty)$.
\end{lemma}
\begin{proof}
\leanok
At each $x \ge 0$ and $r > 0$, for $y > x$ close to $x$: maximizing indices satisfy $f_i(y) < M(x) + r(y-x)$ because
their slopes tend to $f_i'(x) \le 0 < r$, and the other indices stay below $M(x)$ by continuity. So the right slope of $M$
is eventually below $r$. The barrier lemma (Mathlib's \texttt{image\_le\_of\_liminf\_slope\_right\_lt\_deriv\_boundary'})
with $B(t) = M(a) + \varepsilon(t-a)$ gives $M(b) \le M(a) + \varepsilon(b-a)$ for every $\varepsilon > 0$.
\end{proof}

\begin{theorem}[Arrow--Hurwicz: the max-ratio Lyapunov function]
\label{thm:tatonnement_sup}
\lean{Econ.maxRatio, Econ.minRatio, Econ.IsTatonnementPath.hasDerivWithinAt_ratio,
  Econ.tatonnement_maxRatio_antitone, Econ.tatonnement_minRatio_monotone, Econ.tatonnement_bounded}
\leanok
\uses{lem:gs_extreme_ratio, lem:sup_antitone, def:tatonnement}
Let $z$ have gross substitutes and be homogeneous of degree zero, with equilibrium $p^* \gg 0$, for any number of goods.
Along a tâtonnement path that stays in the positive orthant for $t \ge 0$, $\max_i p_i(t)/p^*_i$ is nonincreasing and
$\min_i p_i(t)/p^*_i$ is nondecreasing. Hence $m\,p^* \le p(t) \le M\,p^*$ for all $t \ge 0$, with $m, M$ the extreme
relative prices at time $0$. Invariance of this box implies Lyapunov stability of each equilibrium $c p^*$ (not stated
formally). Convergence is not claimed.
\end{theorem}
\begin{proof}
\leanok
Apply Lemma~\ref{lem:sup_antitone} to $f_i(t) = p_i(t)/p^*_i$, whose derivative is $z_i(p(t))/p^*_i$, nonpositive at a
maximizing index by Lemma~\ref{lem:gs_extreme_ratio}; and to $-f_i$ for the minimum.
\end{proof}

\begin{theorem}[A nonconstant witness]
\label{thm:tatonnement_sup_example}
\lean{Econ.TatonnementSupExamples.zEx, Econ.TatonnementSupExamples.pEx,
  Econ.TatonnementSupExamples.grossSubstitutes_zEx, Econ.TatonnementSupExamples.homogeneousZero_zEx,
  Econ.TatonnementSupExamples.zEx_one, Econ.TatonnementSupExamples.pEx_pos, Econ.TatonnementSupExamples.zEx_pEx,
  Econ.TatonnementSupExamples.isTatonnementPath_pEx, Econ.TatonnementSupExamples.pEx_bounded,
  Econ.TatonnementSupExamples.pEx_nonconstant}
\leanok
\uses{thm:tatonnement_sup}
$z(p) = \bigl((p_1 - p_0)/(p_0+p_1),\ (p_0 - p_1)/(p_0+p_1)\bigr)$ has gross substitutes, is homogeneous of degree zero,
and vanishes at $(1,1)$. The path $p(t) = (1 - e^{-t}/2,\ 1 + e^{-t}/2)$ is a nonconstant tâtonnement path in the
positive orthant, so Theorem~\ref{thm:tatonnement_sup} applies. ($z$ does not satisfy Walras' law, which the theorem
does not use.)
\end{theorem}
\begin{proof}
\leanok
Direct computation: $p_1 - p_0 = e^{-t}$ and $p_0 + p_1 = 2$, so $z(p(t)) = (e^{-t}/2, -e^{-t}/2) = \dot p(t)$.
\end{proof}

\section{Sonnenschein--Mantel--Debreu: Debreu's decomposition}
\label{sec:smd}

\begin{definition}[Open simplex; Debreu's pieces]
\label{def:debreu_piece}
\lean{Econ.openPriceSimplex, Econ.debreuPiece}
\leanok
$\Delta^\circ = \{p \gg 0 : \sum_i p_i = 1\}$. For $z$, a constant $c$ and a good $l$,
$z^l(p) = (z_l(p) + c\, p_l)\,\big(e_l - \tfrac{p_l}{\|p\|^2} p\big)$.
\end{definition}

\begin{lemma}[Each piece satisfies Walras' law; the pieces sum to $z$]
\label{lem:debreu_pieces}
\lean{Econ.cost_self_pos, Econ.cost_single_right, Econ.cost_debreuDir, Econ.walrasLaw_debreuPiece,
  Econ.sum_debreuPiece, Econ.continuousOn_debreuPiece}
\leanok
\uses{def:debreu_piece}
$p \cdot (e_l - \frac{p_l}{\|p\|^2} p) = p_l - p_l = 0$, so $p \cdot z^l(p) = 0$. If $p \cdot z(p) = 0$ then
$\sum_l z^l(p) = z(p)$. Each $z^l$ is continuous wherever $z$ is.
\end{lemma}
\begin{proof}
\leanok
With $\alpha = z(p) + c p$: $\sum_l \alpha_l (e_l - \frac{p_l}{\|p\|^2} p) = \alpha - \frac{\alpha \cdot p}{\|p\|^2} p$, and
$\alpha \cdot p = p \cdot z(p) + c\|p\|^2 = c \|p\|^2$, so the sum is $z(p) + cp - cp = z(p)$. Continuity: sums, products
and the quotient $p_l / \|p\|^2$ with $\|p\|^2 > 0$.
\end{proof}

\begin{lemma}[Strict Cauchy--Schwarz on the open simplex]
\label{lem:sq_cost_lt}
\lean{Econ.sq_cost_lt}
\leanok
For distinct $p, q \in \Delta^\circ$, $(p \cdot q)^2 < \|p\|^2 \|q\|^2$.
\end{lemma}
\begin{proof}
\leanok
Lagrange's identity: $\|p\|^2\|q\|^2 - (p\cdot q)^2 = \tfrac12 \sum_{i,j} (p_i q_j - p_j q_i)^2$. If every term vanished then
$p_i \sum_j q_j = q_i \sum_j p_j$, i.e.\ $p_i = q_i$ for all $i$ (both sum to $1$), contradicting $p \ne q$.
\end{proof}

\begin{lemma}[Each piece satisfies the weak axiom]
\label{lem:debreu_warp}
\lean{Econ.excessWARP_debreuPiece}
\leanok
\uses{lem:debreu_pieces, lem:sq_cost_lt, def:law_of_demand}
If $z_l(p) + c\,p_l > 0$ on $D \subseteq \Delta^\circ$, then $z^l$ satisfies WARP on $D$.
\end{lemma}
\begin{proof}
\leanok
Let $p, q \in D$ with $z^l(p) \ne z^l(q)$ (so $p \ne q$) and $p \cdot z^l(q) \le 0$. Write $N_p = \|p\|^2$, $N_q = \|q\|^2$,
$C = p \cdot q > 0$. Then $p \cdot z^l(q) = \alpha(q)(p_l - q_l C / N_q)$ with $\alpha(q) > 0$, so $p_l \le q_l C/N_q$. Hence
$q \cdot z^l(p) = \alpha(p)(q_l - p_l C/N_p)$ and $p_l C / N_p \le q_l C^2/(N_p N_q) < q_l$ by Lemma~\ref{lem:sq_cost_lt}, so
$q \cdot z^l(p) > 0$.
\end{proof}

\begin{theorem}[Debreu's decomposition; SMD, revealed-preference form]
\label{thm:debreu_decomposition}
\lean{Econ.debreu_decomposition, Econ.debreu_decomposition_of_boundedBelow}
\leanok
\uses{lem:debreu_pieces, lem:debreu_warp}
Let $z$ satisfy Walras' law on $D \subseteq \Delta^\circ$ with $z_l(p) + c\,p_l > 0$ on $D$ for some $c$ (in particular if $z$
is bounded below by $-s$ and all prices in $D$ are at least $\varepsilon > 0$, taking $c = \max(s,0)/\varepsilon$). Then $z$
is a sum of $L$ functions, each satisfying Walras' law and the weak axiom on $D$. That each piece is the excess demand
of a rational consumer (Debreu's second step) is Theorem~\ref{thm:smd_consumers} below.
\end{theorem}
\begin{proof}
\leanok
Take the pieces $z^l$ of Definition~\ref{def:debreu_piece} and apply Lemmas~\ref{lem:debreu_pieces} and
\ref{lem:debreu_warp}. For the bounded-below form, $z_l(p) + c p_l > -s + \max(s,0) \ge 0$ since $c p_l \ge \max(s, 0)$.
\end{proof}

\begin{theorem}[An aggregate violating WARP that decomposes into WARP pieces]
\label{thm:zbad_smd}
\lean{Econ.SMDExamples.fBad, Econ.SMDExamples.zBad, Econ.SMDExamples.DBad, Econ.SMDExamples.DBad_subset,
  Econ.SMDExamples.convex_DBad, Econ.SMDExamples.walrasLaw_zBad, Econ.SMDExamples.not_excessWARP_zBad,
  Econ.SMDExamples.zBad_weight_pos, Econ.SMDExamples.zBad_smd}
\leanok
\uses{thm:debreu_decomposition, thm:warp_convex}
With $f(t) = (t - \frac14)(t - \frac34)$ and $z(p) = (f(p_0), -p_0 f(p_0)/p_1)$ on
$D = \{p \in \Delta^\circ : p_i \ge \frac18\}$, $z$ satisfies Walras' law, has equilibria $(\frac14, \frac34)$ and
$(\frac34, \frac14)$ but not their midpoint, hence violates WARP; yet it is a sum of two WARP excess demands.
\end{theorem}
\begin{proof}
\leanok
Walras' law: $p_0 f - p_1 \cdot p_0 f/p_1 = 0$. $D$ is convex. If $z$ satisfied WARP, its equilibrium set would be convex
(Theorem~\ref{thm:warp_convex}), but $f(\frac12) = -\frac1{16} \ne 0$. Weights with $c = 16$: $f(p_0) + 16 p_0 \ge -\frac1{16} + 2$,
and $-p_0 f(p_0)/p_1 + 16 p_1 > 0$ because $p_0 f(p_0) \le \frac3{16} < \frac14 \le 16 p_1^2$. Apply
Theorem~\ref{thm:debreu_decomposition}.
\end{proof}

\subsection{Step 2: each piece is a rational consumer}
\label{sec:smd_consumers}

Fix a good $l$, a compact set $K \subseteq \Delta^\circ$ of prices, and a weight $a$ continuous and positive on $K$. Write
$\|p\| = \sqrt{p\cdot p}$, $V(q) = q_l / \|q\|$ (the cosine of $q$ with the $l$-th axis) and
$y(q) = a(q)\,(e_l - (q_l/\|q\|^2)\,q)$, so Debreu's piece $z^l$ is $y$ with $a(q) = z_l(q) + c\,q_l$.

\begin{definition}[Level sets and utility]
\label{def:piece_utility}
\lean{Econ.dirCos, Econ.weightedPiece, Econ.debreuPiece_eq_weightedPiece, Econ.pieceLevelSet, Econ.pieceUtility}
\leanok
\uses{def:debreu_piece}
For $\theta \in \mathbb{R}$ let $C_\theta = \operatorname{conv}\{\,y(q) : q \in K,\ V(q) \le \theta\,\}$, and for a net trade
$x$ let $u(x) = -\inf\bigl(\{2\} \cup \{\theta : x \in C_\theta\}\bigr)$.
\end{definition}

\begin{lemma}[Key identity]
\label{lem:piece_identity}
\lean{Econ.cost_weightedPiece, Econ.piece_identity, Econ.cs_gap_pos, Econ.cs_gap_nonneg, Econ.dirCos_pos,
  Econ.dirCos_le_one, Econ.cost_comm}
\leanok
\uses{def:piece_utility, lem:sq_cost_lt}
For positive $p, q$: $p \cdot y(q) = a(q)\,\Phi(p,q)$ with
\[
  \Phi(p,q) = p_l - \frac{q_l\,(p\cdot q)}{\|q\|^2}
  = \|p\|\,\bigl(V(p) - V(q)\bigr) + \frac{q_l}{\|q\|^2}\bigl(\|p\|\,\|q\| - p\cdot q\bigr).
\]
On $\Delta^\circ$ the last bracket is $\ge 0$, and $> 0$ for $q \ne p$. Also $0 < V(q) \le 1$.
\end{lemma}
\begin{proof}
\leanok
Expand $y(q)$ and use $p \cdot e_l = p_l$; the identity is algebra after writing $\|q\|^2 = \sqrt{q\cdot q}^{\,2}$. The bracket is
positive by strict Cauchy--Schwarz (Lemma~\ref{lem:sq_cost_lt}): $(p\cdot q)^2 < \|p\|^2\|q\|^2$ and $p \cdot q > 0$, so
$p\cdot q < \|p\|\|q\|$. $V(q) \le 1$ because $q_l^2 \le \sum_i q_i^2$.
\end{proof}

\begin{lemma}[Strict separation below the level of $p$]
\label{lem:piece_sep}
\lean{Econ.exists_comb_of_mem_pieceLevelSet, Econ.cost_comb, Econ.wsum_pos, Econ.cost_pos_of_mem_pieceLevelSet}
\leanok
\uses{lem:piece_identity}
If $p \in K$ and $\theta < V(p)$, every $x \in C_\theta$ has $p \cdot x > 0$.
\end{lemma}
\begin{proof}
\leanok
Write $x = \sum_i w_i\, y(q_i)$ with $w_i \ge 0$, $\sum w_i = 1$, $V(q_i) \le \theta$. By Lemma~\ref{lem:piece_identity},
$\Phi(p, q_i) \ge \|p\|(V(p) - \theta) > 0$, so $p \cdot x = \sum_i w_i\, a(q_i)\,\Phi(p,q_i) > 0$ (some $w_i > 0$).
\end{proof}

\begin{lemma}[Quantitative uniqueness]
\label{lem:piece_gap}
\lean{Econ.exists_gap, Econ.continuousOn_weightedPiece}
\leanok
\uses{lem:piece_identity}
Let $p \in K$ and $x \ne y(p)$ with $p \cdot x \le 0$. There is $\delta > 0$ such that $x \notin C_\theta$ for all
$\theta < V(p) + \delta$.
\end{lemma}
\begin{proof}
\leanok
Let $\eta = \|x - y(p)\|/2$ and $F = \{q \in K : \|y(q) - y(p)\| \ge \eta\}$, compact and not containing $p$. Let
$g(q) = \frac{q_l}{\|q\|^2}(\|p\|\|q\| - p\cdot q)$; by compactness $g \ge \mu > 0$ on $F$, $m \le a \le M$ on $K$ with $m > 0$,
and $\|y(q) - y(p)\| \le B$ on $K$. Put $D = 2BM\|p\|/(m\mu)$ and $\delta = \min\bigl(\mu/(2\|p\|),\, \eta/(D+1)\bigr)$.

Suppose $x = \sum_i w_i y(q_i) \in C_\theta$ with $\theta < V(p) + \delta$. By the identity, $\Phi_i := \Phi(p,q_i) > -\|p\|\delta + g(q_i)$.
Let $r_i = a(q_i)\Phi_i + M\|p\|\delta \ge 0$. If $q_i \in F$ then $\Phi_i \ge \mu/2$, so $r_i \ge m\mu/2$; hence for every $i$,
$\|y(q_i) - y(p)\| \le \eta + \frac{2B}{m\mu} r_i$. Since $\sum_i w_i r_i = p \cdot x + M\|p\|\delta \le M\|p\|\delta$,
\[
  \|x - y(p)\| \le \sum_i w_i \|y(q_i) - y(p)\| \le \eta + D\delta < 2\eta = \|x - y(p)\|,
\]
a contradiction.
\end{proof}

\begin{theorem}[The piece is a convex consumer's unique choice]
\label{thm:piece_utility}
\lean{Econ.bddBelow_levels, Econ.pieceLevelSet_mono, Econ.quasiconcaveOn_pieceUtility, Econ.pieceUtility_self,
  Econ.pieceUtility_lt}
\leanok
\uses{lem:piece_sep, lem:piece_gap}
$u$ is quasiconcave, $u(y(p)) = -V(p)$ for $p \in K$, and $u(x) < u(y(p))$ for every $x \ne y(p)$ with $p\cdot x \le 0$.
\end{theorem}
\begin{proof}
\leanok
The level sets of $\theta$ are bounded below by $0$ (as $V > 0$) and increase with $\theta$. Quasiconcavity: if
$\inf\Theta(x_1), \inf\Theta(x_2) \le t$ then for each $\varepsilon > 0$ pick $\theta_j < t + \varepsilon$ in the two sets; either one is
the floor $2$, or $x_1, x_2 \in C_{\max\theta_j}$, which is convex, so the mixture is in it. $u(y(p)) = -V(p)$: $y(p) \in C_{V(p)}$,
and $y(p) \notin C_\theta$ for $\theta < V(p)$ by Lemma~\ref{lem:piece_sep} since $p \cdot y(p) = 0$. Strictness:
Lemma~\ref{lem:piece_gap} gives $\inf\Theta(x) \ge V(p) + \min(\delta, 1)$ (the floor $2 \ge V(p) + 1$).
\end{proof}

\begin{theorem}[A rational consumer for each piece]
\label{thm:consumer_piece}
\lean{Econ.exists_consumer_weightedPiece}
\leanok
\uses{thm:piece_utility}
There are a quasiconcave $U$ and an endowment $\omega \gg 0$ such that for all $p \in K$ the Walrasian demand of $U$ at
income $p \cdot \omega$ is exactly $\{\omega + y(p)\}$.
\end{theorem}
\begin{proof}
\leanok
Let $b$ bound $\|y\|$ on $K$, $\omega = \max(b, 1)\mathbf{1}$, and $U(x) = u(x - \omega)$ (a translate, so quasiconcave). Then
$\omega + y(p) \ge 0$ and costs exactly $p \cdot \omega$. Any other affordable $x \ge 0$ has $p\cdot(x - \omega) \le 0$ and
$x - \omega \ne y(p)$, so $U(x) < U(\omega + y(p))$ by Theorem~\ref{thm:piece_utility}.
\end{proof}

\begin{theorem}[Sonnenschein--Mantel--Debreu, with convex consumers]
\label{thm:smd_consumers}
\lean{Econ.sonnenschein_mantel_debreu_quasiconcave}
\leanok
\uses{thm:consumer_piece, thm:debreu_decomposition}
Let $z$ be continuous and satisfy Walras' law on a compact $K \subseteq \Delta^\circ$. Then there are $L$ consumers with
quasiconcave utilities $U_l$ and endowments $\omega_l \gg 0$ whose Walrasian demands at every $p \in K$ are single bundles
$x_l(p)$ with $\sum_l (x_l(p) - \omega_l) = z(p)$. Caveat: Debreu's consumers also have continuous, strictly convex, monotone
preferences. The utilities here provably lack all three ($u = -2$ off the bounded set $C_1$, so they jump and are satiated);
standard consumers need a different construction (Geanakoplos 1984), which is open.
\end{theorem}
\begin{proof}
\leanok
$q \mapsto z_l(q)/q_l$ is continuous on the compact $K$, hence $\ge b_l$; with $c = 1 + \sum_l |b_l|$ the weight
$z_l(q) + c q_l \ge q_l > 0$. Apply Theorem~\ref{thm:consumer_piece} to each Debreu piece and sum
(Lemma~\ref{lem:debreu_pieces}).
\end{proof}

\begin{theorem}[A WARP-violating aggregate from two rational consumers]
\label{thm:zbad_consumers}
\lean{Econ.SMDExamples.isCompact_DBad, Econ.SMDExamples.continuousOn_zBad,
  Econ.SMDExamples.zBad_rational_consumers}
\leanok
\uses{thm:smd_consumers, thm:zbad_smd}
The excess demand $z$ of Theorem~\ref{thm:zbad_smd} violates WARP, yet is the aggregate excess demand of two consumers as in
Theorem~\ref{thm:smd_consumers} on $D$.
\end{theorem}
\begin{proof}
\leanok
$D$ is closed (an equality and inequalities) and inside $[0,1]^2$, hence compact; $z$ is continuous on $D$ since $p_1 \ge \frac18$.
\end{proof}

\subsection{Step 2 with continuous, strongly monotone consumers}
\label{sec:smd_standard}

Keep the notation of \S\ref{sec:smd_consumers}, write $d(q) = e_l - (q_l/\|q\|^2)q$ (so $y(q) = a(q)d(q)$), and assume all
prices in $K$ are at least $2\varepsilon > 0$. Let $Q = \{q : \varepsilon \le q_i \le 2\}$.

\begin{definition}[Level test and utility]
\label{def:std_utility}
\lean{Econ.cnorm, Econ.dvec, Econ.csGap, Econ.priceBox, Econ.levelGap, Econ.ReachesLevel, Econ.stdUtility,
  Econ.PieceData, Econ.weightedPiece_eq_smul_dvec}
\leanok
\uses{def:piece_utility}
For $q \in Q$, $p' \in K$ and $\theta \in \mathbb{R}$ let $G(q,p',\theta) = a(p')\bigl(q\cdot d(p') + \|q\|(V(p') - \theta)\bigr)$.
A net trade $x$ \emph{reaches} $\theta$ if for every $q \in Q$ some $p' \in K$ has $G(q,p',\theta) \le q\cdot x$. Let
$u(x) = -\inf\{\theta : x \text{ reaches } \theta\}$.
\end{definition}

\begin{lemma}[Euclidean estimates]
\label{lem:std_estimates}
\lean{Econ.cost_self_nonneg', Econ.cnorm_nonneg, Econ.cnorm_sq, Econ.cnorm_mul_self, Econ.cost_sq_le,
  Econ.abs_cost_le, Econ.cost_le_cnorm_mul, Econ.cost_sub_sub, Econ.cost_comm', Econ.abs_cnorm_sub_le,
  Econ.abs_apply_le_cnorm, Econ.cnorm_mul_sub_cost_le, Econ.cost_dvec, Econ.csGap_nonneg, Econ.cost_dvec_eq,
  Econ.dirCos_eq, Econ.cnorm_pos, Econ.abs_dirCos_sub_le, Econ.abs_csGap_sub_le, Econ.cnorm_ge_of_mem_priceBox,
  Econ.cnorm_le_of_mem_priceBox, Econ.abs_cost_le_of_mem_priceBox, Econ.abs_cost_dvec_add_le,
  Econ.sInf_le_sInf_add, Econ.exists_eps_of_compact}
\leanok
\uses{lem:piece_identity}
Cauchy--Schwarz $|p\cdot q| \le \|p\|\|q\|$ and $\bigl|\|q\| - \|p\|\bigr| \le \|q - p\|$. Also
$\|q\|\|p\| - q\cdot p \le \|q-p\|^2/2$ and $|V(q) - V(p)| \le 2\|q-p\|/\|q\|$. The gap
$g(q,p') = \frac{p'_l}{\|p'\|^2}(\|q\|\|p'\| - q\cdot p') \ge 0$ satisfies $|g(q,p') - g(p,p')| \le 2\|q - p\|$. On $Q$:
$\varepsilon \le \|q\| \le 2L$, $|q \cdot x| \le 2L\|x\|_\infty$ and $|q\cdot d(p') + \|q\|V(p')| \le 6L$. Finally,
$q \cdot d(p') = \|q\|(V(q) - V(p')) + g(q,p')$.
\end{lemma}
\begin{proof}
\leanok
Cauchy--Schwarz by Lagrange's identity. The others follow by expanding $\|q - p\|^2 = \|q\|^2 + \|p\|^2 - 2q\cdot p$ and using
$|q_i| \le \|q\|$. The last identity is Lemma~\ref{lem:piece_identity}.
\end{proof}

\begin{theorem}[The utility is continuous, strongly monotone and quasiconcave]
\label{thm:std_utility_props}
\lean{Econ.PieceData.nonemptyι, Econ.PieceData.mem_priceBox, Econ.PieceData.exists_min, Econ.PieceData.exists_max,
  Econ.PieceData.levels_nonempty, Econ.PieceData.levels_bddBelow, Econ.PieceData.reaches_mono,
  Econ.PieceData.reaches_shift, Econ.PieceData.stdUtility_lipschitz, Econ.PieceData.continuous_stdUtility,
  Econ.PieceData.stdUtility_strictMono, Econ.PieceData.quasiconcaveOn_stdUtility, Econ.levelGap_add}
\leanok
\uses{def:std_utility, lem:std_estimates}
$u$ is Lipschitz with constant $2L/(m\varepsilon)$, strictly increasing in the coordinatewise order, and quasiconcave.
\end{theorem}
\begin{proof}
\leanok
$G(q,p',\theta+\Delta) = G(q,p',\theta) - a(p')\|q\|\Delta$ with $m\varepsilon \le a\|q\| \le 2LM$. The set of reached
levels is nonempty and bounded below by the bounds on $Q$. If $x$ reaches $\theta$, then $x'$ reaches
$\theta + 2L\|x - x'\|/(m\varepsilon)$ because $|q\cdot(x - x')| \le 2L\|x - x'\|$, so $u$ is Lipschitz. If $x \le x'$ with
$x_i < x'_i$, then $q \cdot (x' - x) \ge \varepsilon(x'_i - x_i)$ on $Q$. So $x'$ reaches $\theta - \varepsilon(x'_i - x_i)/(2LM)$,
and $u(x') > u(x)$. Each test is a half-space in $x$: given $p_1, p_2$ for $x_1, x_2$, the one with smaller $G$ works for any
convex combination. Hence the reached sets are convex and increase with $\theta$, which gives quasiconcavity as in
Theorem~\ref{thm:piece_utility}.
\end{proof}

\begin{lemma}[Quantitative uniqueness]
\label{lem:std_gap}
\lean{Econ.reaches_self, Econ.PieceData.not_reaches_below, Econ.PieceData.exists_gap,
  Econ.PieceData.stdUtility_self, Econ.PieceData.stdUtility_lt}
\leanok
\uses{thm:std_utility_props}
$u(y(p)) = -V(p)$ for $p \in K$, and $u(x) < u(y(p))$ whenever $p \cdot x \le 0$ and $x \ne y(p)$.
\end{lemma}
\begin{proof}
\leanok
$y(p)$ reaches $V(p)$ (take $p' = p$: $G(q,p,V(p)) = q\cdot y(p)$). At $q = p$, $G(p,p',\theta) = a(p')(g(p,p') + \|p\|(V(p) - \theta)) > 0$
for $\theta < V(p)$, so nothing with $p\cdot x \le 0$ reaches below $V(p)$.

For $x \ne y(p)$ with $p\cdot x \le 0$, let $w = x - y(p)$, $\gamma = \|w\|^2 > 0$ and test with $q = p - t w$ for small $t > 0$.
Then $q \in Q$, $\|q - p\| = t\|w\|$ and $q \cdot x \le q\cdot y(p) - t\gamma$. Write $E(q,p') = q\cdot d(p') + \|q\|(V(p') - V(p)) =
D + g(q,p')$ with $D = \|q\|(V(q) - V(p))$, $|D| \le N C_V t$, and $g(q,p) \le C_g t^2$. Then
$a(p')E(q,p') - a(p)E(q,p) = (a(p') - a(p))D + a(p')g(q,p') - a(p)g(q,p)$.

Choose $\eta_0 = \gamma/(4NC_V)$. When $|a(p') - a(p)| < \eta_0$ this is $\ge -\gamma t/2$. Otherwise $p'$ lies in a compact set
avoiding $p$, where $g(p,\cdot) \ge \mu > 0$. Then $a(p')g(q,p') \ge m(\mu - 2t\|w\|)$ dominates for small $t$. Since
$a(p)E(q,p) = q\cdot y(p)$, every $p'$ gives
$G(q,p',\theta) \ge q\cdot y(p) - \tfrac34\gamma t > q\cdot x$ for $\theta < V(p) + \delta$ with $\delta = \gamma t/(4MN)$.
So $x$ does not reach any such $\theta$, and $u(x) \le -V(p) - \delta$.
\end{proof}

\begin{theorem}[A standard consumer for each piece]
\label{thm:std_consumer_piece}
\lean{Econ.exists_standard_consumer_weightedPiece}
\leanok
\uses{lem:std_gap}
There are a continuous, strongly monotone, quasiconcave $U$ and an endowment $\omega \gg 0$ such that, for all $p \in K$, the
Walrasian demand of $U$ at income $p\cdot\omega$ is exactly $\{\omega + y(p)\}$.
\end{theorem}
\begin{proof}
\leanok
$\varepsilon$ exists by compactness. Take $U(x) = u(x - \omega)$ with $\omega = \max(b,1)\mathbf 1$, where $b$ bounds $\|y\|$ on $K$, and
argue as in Theorem~\ref{thm:consumer_piece}. If $K$ is empty, $U(x) = \sum_i x_i$ works.
\end{proof}

\begin{theorem}[Sonnenschein--Mantel--Debreu, with standard consumers]
\label{thm:smd_standard}
\lean{Econ.sonnenschein_mantel_debreu_monotone, Econ.SMDExamples.zBad_standard_consumers}
\leanok
\uses{thm:std_consumer_piece, thm:smd_consumers}
Let $z$ be continuous and satisfy Walras' law on a compact $K \subseteq \Delta^\circ$. Then $z$ is the aggregate excess demand on $K$
of $L$ consumers with continuous, strongly monotone, quasiconcave utilities and endowments $\omega_l \gg 0$, each with single-valued
demand. In particular, the WARP-violating $z$ of Theorem~\ref{thm:zbad_smd} arises this way. Caveat: Debreu's preferences are also
strictly convex. These are only convex, so individual demand is guaranteed single-valued only on $K$.
\end{theorem}
\begin{proof}
\leanok
As in Theorem~\ref{thm:smd_consumers}, using Theorem~\ref{thm:std_consumer_piece} for each piece.
\end{proof}
